\chapter{at-least-once에서 exactly-once로}
\label{ch:delivery}

\paragraph{In brief.}
\begin{description}
  \item[Why] 스트리밍 파이프라인은 "죽었다 다시 살아나는" 일이 늘 일어난다. 그때마다
    메시지가 사라지거나(유실) 두 번 들어오는(중복) 일이 생기는데, 이 파이프라인이
    감당해야 하는 진짜 문제는 유실을 완전히 없애는 것이 아니라 \term{어디서 유실이
    나는지}를 정확히 알고, \term{그 나머지 구간}에서는 중복을 안전하게 접는 것이다.
  \item[When] producer가 거래소 WS(웹소켓, WebSocket) 연결을 받아 Kafka로 내보내는
    순간부터, ClickHouse의 데이터 테이블에 행이 최종적으로 쌓이는 순간까지 전 구간에
    적용된다.
  \item[Where] 이 장의 결론은 예측시장(Kalshi, Polymarket)과 현물(Binance, Coinbase)
    가족 모두에 동일하게 적용된다. \texttt{venue} 축(\Cref{ch:contract})이 이 구간의
    보장 수준을 바꾸지 않기 때문이다.
  \item[How] 앱 내부 메모리 큐 구간은 보장이 아예 없다는 사실을 받아들이고, 그 뒤
    구간(브로커부터 ClickHouse까지)은 \term{정렬 키(sort key)가 완전히 같은 행만 접는다}는
    ClickHouse \texttt{ReplacingMergeTree}의 규칙 하나로 중복을 흡수한다.
\end{description}

\section{배경: "적어도 한 번"과 "정확히 한 번"이라는 말이 왜 미끄러운가}
\label{sec:delivery-vocab}

스트리밍 파이프라인을 처음 설계할 때 누구나 듣는 말이 있다. "at-least-once(적어도
한 번 전달)로 받고 exactly-once(정확히 한 번)로 저장한다." 문제는 이 두 단어가
\term{파이프라인 전체}의 성질처럼 들리지만, 실제로는 \term{구간마다 다르게 성립}한다는
점이다. 한 파이프라인 안에 유실이 나는 구간과 중복만 나는 구간이 동시에 있을 수 있고,
"이 시스템은 at-least-once다"라는 한 문장은 그 차이를 지워 버린다.

이 파이프라인은 producer(거래소 WS를 구독해 Kafka로 내보내는 프로세스) 3종에서
시작해 Kafka, 그리고 ClickHouse의 데이터 테이블까지 이어지는 사슬이다. 이 장은 그
사슬을 네 구간으로 쪼개, 각 구간이 실제로 어떤 보장을 갖는지 하나씩 확인한다.

\term{at-least-once(적어도 한 번 전달)}란 "메시지가 유실되지는 않지만, 같은 메시지가
두 번 이상 도착할 수 있다"는 보장이다. 반대로 \term{exactly-once(정확히 한 번)}는
"유실도 없고 중복도 없다"는, 더 강하지만 훨씬 지키기 어려운 보장이다. 스트리밍
시스템에서 진짜 exactly-once를 배달 단계에서 보장하는 방법은 사실상 없다시피 하다.
대신 이 파이프라인이 택한 길은 "배달은 중복을 허용하되, 저장 단계에서 중복을
접어 \term{effectively-exactly-once(사실상 정확히 한 번)}를 만든다"이다.
\Cref{sec:delivery-dedup}에서 그 방법을 다룬다.

\section{앱 내부 큐: at-least-once조차 아닌 구간}
\label{sec:delivery-leak}

\paragraph{무슨 문제가 있었나.}
producer가 거래소 WS로부터 메시지를 받으면, 그 메시지를 곧바로 Kafka로 보내지 않는다.
먼저 애플리케이션 안의 큐(\texttt{BufferedKafkaProducer.queue})에 쌓아 두고,
\texttt{BATCH\_SIZE}(100개) 또는 \texttt{FLUSH\_MS}(2000ms, 2초) 조건 중 먼저 차는
쪽에서 한꺼번에 드레인(비워서 내보냄)한다.
% src: 00.architecture.mvp1.md §5.3-1

이 큐는 \term{순수 메모리(in-memory)}다. 디스크에 쓰지 않고, 다른 프로세스로 복제하지도
않는다. 그래서 producer 프로세스가 죽으면, 그 순간 큐 안에 있던 메시지는 그대로
사라진다. 최악의 경우 \texttt{FLUSH\_MS}가 다 차기 직전, 즉 최대 2초분의 메시지가
한 번에 사라질 수 있다.
% src: 00.architecture.mvp1.md §5.3-1

\begin{callout}[왜 이것이 "조용한" 유실인가]
여기서 진짜 문제는 유실이 일어난다는 사실 자체가 아니다. \term{그 유실이 아무 흔적도
남기지 않는다}는 점이다. 두 가지가 겹쳐서 이런 결과를 만든다.

\begin{enumerate}
  \item \texttt{dr\_cb: false}. \term{delivery report 콜백}은 Kafka 클라이언트
    라이브러리(librdkafka)가 "이 메시지가 결국 브로커에 쓰였는지, 아니면 끝내
    실패했는지"를 앱에 알려 주는 통로다. 이 콜백이 꺼져 있으면, produce가 최종적으로
    실패해도 애플리케이션은 그 사실을 알 방법이 없다.
    % src: 00.architecture.mvp1.md §5.3-2
  \item WS 피드는 \term{재생(replay)이 없다}. 거래소가 "5초 전에 보낸 그 메시지를
    다시 보내 달라"는 요청을 받아 주지 않는다. 한 번 놓친 메시지는 그 순간의 호가·체결
    정보를 영원히 잃는 것과 같다.
\end{enumerate}
\end{callout}

\paragraph{왜 이것이 "최악의 종류"의 버그인가.} 버그에는 두 종류가 있다. 하나는
터지는 순간 로그나 에러 메시지, 크래시로 자기 존재를 알리는 버그다. 다른 하나는
아무 신호도 남기지 않고 그냥 결과가 조용히 틀려지는 버그다. 이 큐 유실은 후자다.
프로세스는 정상 종료 로그를 남기고, 재시작 후에도 아무 에러 없이 다시 동작한다.
잃어버린 2초분의 데이터는 로그 어디에도, 메트릭 어디에도 나타나지 않는다. 유실이
있었다는 사실을 알아채려면 ClickHouse에 쌓인 시계열에서 \term{시간 구멍}을 직접
찾아내는 수밖에 없는데, BTC 1종처럼 메시지가 촘촘한 스트림에서는 그 구멍조차
눈에 띄기 어렵다. 발견되지 않는 버그는 고쳐지지도 않는다. 이것이 "발생을 완전히
막을 수는 없어도, 최소한 발생했다는 사실만은 반드시 알 수 있어야 한다"는 원칙이
스트리밍 파이프라인 설계에서 특히 중요한 이유다.

여기에 셧다운 경합(race) 하나가 더 있다. 프로세스가 종료 신호를 받은 뒤에도 아직
처리 중이던 메시지가 큐에 넣으려 시도하면, 이미 닫힌 연결이라 그 메시지는 버려지고
\texttt{warn} 로그 한 줄만 남는다(\texttt{enqueue after disconnect dropped}). 이
경우는 앞의 두 경우와 달리 로그가 남기는 하지만, error가 아니라 warn이라 운영
관점에서 놓치기 쉽다.
% src: 00.architecture.mvp1.md §5.3-3

정리하면, producer 내부 큐부터 Kafka 브로커에 데이터가 실제로 쓰이기까지의 구간은
\term{at-least-once조차 성립하지 않는다}. "적어도 한 번은 전달된다"는 말 자체가
이 구간에는 적용되지 않으며, 그 사실이 로그에도 남지 않는다.

\section{브로커부터 ClickHouse까지: at-least-once가 성립하는 구간}
\label{sec:delivery-atleastonce}

\paragraph{무슨 문제가 있었나.} 메시지가 큐를 벗어나 Kafka 클라이언트 내부 버퍼로
넘어간 다음부터는 사정이 다르다. 여기서부터는 네트워크 문제나 브로커 일시 장애로
인한 실패를 \term{재시도(retry)}로 극복할 수 있는 구간이기 때문이다. 문제는 "재시도가
있다"는 사실만으로는 안심할 수 없다는 것이다. 재시도가 있어도 ack 수준이 약하면
여전히 유실이 날 수 있고, 재시도가 만드는 중복을 받아줄 준비가 없으면 중복 자체가
새로운 문제가 된다.

\paragraph{성립하는 이유.} librdkafka(Kafka 클라이언트 라이브러리)는 ack(승인)를 받지
못한 produce 요청을 \texttt{message.timeout.ms}까지 자동으로 재시도한다. 이 재시도는
\term{같은 바이트}를 다시 보내는 것이므로, 같은 \texttt{idempotency\_key}(멱등성 키)와
같은 \texttt{time\_ms}를 가진 메시지가 그대로 다시 전송된다. 즉 여기서 생기는 중복은
전부 \term{완전히 동일한 행}이다. 이 성질이 \Cref{sec:delivery-dedup}의 중복 제거를
가능하게 만드는 전제가 된다.
% src: 00.architecture.mvp1.md §5.2 (2)

ClickHouse의 \texttt{ENGINE=Kafka} 테이블도 같은 원리로 동작한다. Kafka 엔진은
메시지를 읽어 \term{materialized view(구체화 뷰, MV)}에 삽입을 마친 \term{후에만}
그 오프셋(offset, 읽은 위치)을 커밋한다. 삽입과 커밋 사이에서 ClickHouse가 죽으면,
재시작 후 같은 오프셋부터 다시 읽어 같은 블록을 다시 삽입한다. 이 역시 완전히 동일한
행의 중복을 만든다.
% src: 00.architecture.mvp1.md §5.2 (3)

\paragraph{단서 하나: ack 수준.} 코드는 \texttt{acks} 옵션을 명시하지 않으므로
librdkafka의 기본값인 \texttt{request.required.acks = -1}(all, 모든 복제본의 ack를
기다림)이 적용된다. 그런데 이 파이프라인의 Kafka 클러스터는 \term{replication
factor(복제 계수)가 1}이다. 복제본이 하나뿐이면 "모든 복제본의 ack"는 결국 "리더
브로커 한 대의 ack"와 같은 뜻이 된다. 즉 설정값은 \texttt{acks=all}이지만 \term{실질
효과는 \texttt{acks=1}}이다. 그 한 대가 디스크에 fsync(강제로 디스크에 씀)하기 전에
죽으면, 재시도로도 메꿀 수 없는 유실이 난다. 이 단서는 \Cref{sec:delivery-honest}에서
"정직하게 어디까지가 남은 위험인가"를 정리할 때 다시 등장한다.
% src: 00.architecture.mvp1.md §5.2 (2)

정리하면, \term{브로커부터 ClickHouse 데이터 테이블까지는 at-least-once가 성립한다.}
재시도와 지연 커밋이 만드는 중복은 전부 바이트 단위로 동일한 재전송이라는 성질을
가지며, 이 성질을 \Cref{sec:delivery-dedup}이 이용한다.

\section{파이프라인 전체 지도}
\label{sec:delivery-map}

앞 두 절의 내용을 한 장의 그림으로 모으면 \Cref{fig:delivery-stages}와 같다. 화살표
왼쪽의 괄호 번호는 이 장에서 다룬 구간 번호이고, 화살표 위 라벨이 그 구간의 실제
보장 수준이다.

\begin{figure}[htbp]
  \centering
  \begin{tikzpicture}[
      node distance=8mm and 14mm,
      box/.style={draw, rounded corners, minimum width=32mm, minimum height=9mm,
                   align=center, font=\small},
      lab/.style={font=\scriptsize, align=center},
      arr/.style={-{Stealth[length=2.2mm]}, thick}
    ]
    \node[box] (ws) {거래소 WS};
    \node[box, right=of ws] (queue) {producer\\메모리 큐};
    \node[box, right=of queue] (broker) {Kafka 브로커\\(파티션 1, RF 1)};
    \node[box, right=of broker] (mv) {ClickHouse\\Kafka 엔진 $\to$ MV};
    \node[box, right=of mv] (table) {ReplacingMergeTree\\데이터 테이블};

    \draw[arr, red!60!black] (ws) -- node[lab, above]{(1) 보장 없음\\최대 2초분 유실 가능} (queue);
    \draw[arr] (queue) -- node[lab, above]{(2) at-least-once\\재시도, 동일 바이트 재전송} (broker);
    \draw[arr] (broker) -- node[lab, above]{(3) at-least-once\\MV 삽입 후 offset 커밋} (mv);
    \draw[arr, blue!50!black] (mv) -- node[lab, above]{(4) sort key 완전 일치 시\\병합에서 중복 제거} (table);
  \end{tikzpicture}
  \caption{producer 큐부터 ClickHouse 데이터 테이블까지, 구간별 전달 보장 수준.
    (1)만 보장이 없고 나머지는 at-least-once이며, (4)에서 effectively-exactly-once로
    접힌다.}
  \label{fig:delivery-stages}
\end{figure}
% src: 00.architecture.mvp1.md §5.1

\section{DB 단에서 exactly-once로 접는 방법}
\label{sec:delivery-dedup}

\paragraph{무슨 문제가 있었나.} \Cref{sec:delivery-atleastonce}에서 확인했듯, 브로커
이후 구간은 중복을 만들 수는 있어도 그 중복은 전부 바이트 단위로 동일한 재전송이다.
그렇다면 남은 문제는 하나다. \term{동일한 행이 여러 번 들어왔을 때, 저장 단계에서
그것을 자동으로 한 행으로 되돌릴 방법}이 필요하다는 것이다.

\paragraph{그래서 무엇을 바꿨나.} 데이터 테이블의 엔진으로
\texttt{ReplacingMergeTree(ingested\_at)}를 쓰고, 정렬 키(sort key)를
\texttt{ORDER BY (venue, ticker, time, idempotency\_key)}로 잡는다.
% src: 00.architecture.mvp1.md §3.4, §5.4

\paragraph{왜 그 변경이 문제를 푸는가.} \term{ReplacingMergeTree}는 백그라운드에서
파트(part, 저장 단위)를 병합(merge)할 때, \term{정렬 키가 완전히 같은 행들만} 하나로
접는다. 여기서 "완전히 같다"는 것은 \texttt{idempotency\_key} 하나만 같아서는 안 되고,
\texttt{venue}, \texttt{ticker}, \texttt{time}까지 \term{네 컬럼 전부}가 같아야
한다는 뜻이다. 정렬 키의 일부만 같은 행은 별개의 행으로 남는다.
% src: 00.architecture.mvp1.md §5.4

이 규칙이 이 파이프라인에서 실제로 작동하는 이유는, \Cref{sec:delivery-atleastonce}에서
확인한 성질 때문이다. \term{이 파이프라인의 모든 중복은 바이트 단위로 동일한
재전송}이므로, 같은 이벤트가 두 번 들어오면 \texttt{venue}, \texttt{ticker},
\texttt{time}, \texttt{idempotency\_key} 네 값이 매번 똑같다. 여기서 \texttt{time}은
ClickHouse가 메시지를 받은 시각이 아니라, MV가 페이로드의 \texttt{time\_ms}를
\texttt{toDateTime64(time\_ms/1000, 3)}으로 변환해 얻은 값이다. 즉 \term{payload
안의 값에서 나온다}. 재전송이 일어나도 페이로드 자체는 바뀌지 않으므로 \texttt{time}도
매번 같고, 결과적으로 정렬 키 네 컬럼 전부가 자동으로 일치해 병합 시 접힌다.
% src: 00.architecture.mvp1.md §5.4

\begin{callout}[이 불변식을 깨는 유일한 경우]
정렬 키가 완전히 같아야 접힌다는 규칙은 거꾸로, \term{애플리케이션이 같은 이벤트를
다른 \texttt{time\_ms}로 다시 만들어 내면} 중복 제거가 작동하지 않는다는 뜻이기도
하다. 예를 들어 producer가 재시작한 뒤, 재시도로 보낸 것이 아니라 \term{같은 오더북
상태를 새로운 \texttt{Date.now()}로 다시 publish}하면, 그 두 행은 \texttt{time}이
달라 정렬 키가 일치하지 않고 이력에 별개의 두 행으로 남는다. 이것은 재전송이 아니라
\term{새로운 이벤트로 취급되는 것}이므로, ReplacingMergeTree 입장에서는 정상 동작이다.
% src: 00.architecture.mvp1.md §5.4
\end{callout}

\paragraph{\texttt{venue}가 정렬 키 선두에 있어야 하는 이유.} Kalshi의 티커는
\texttt{KXBTCD-26SEP22-B112000}처럼 시리즈 형태 문자열이고, Polymarket의 티커는
CLOB(Central Limit Order Book) token id로, 수십 자리 10진수 문자열이다. 두 거래소의
티커 값이 \term{우연히 같은 문자열}이 되는 일은 실무적으로 없다고 봐도 되지만, 그
가능성을 완전히 배제하려면 \texttt{venue}가 정렬 키와 \texttt{idempotency\_key} 양쪽
모두에 들어가야 한다.
% src: 00.architecture.mvp1.md §3.4, §5.4

\texttt{venue}가 없다면, 두 거래소의 행이 \texttt{ticker}만으로 구분되므로, 티커
문자열이 우연히 겹치는 극단적인 경우 서로 다른 거래소의 행이 하나로 잘못 접힐 수
있다. \texttt{venue}를 정렬 키 선두에 두면 이 위험이 원천적으로 사라진다. 동시에
\texttt{venue}는 카디널리티(값의 종류 수)가 2~4개로 낮아, ClickHouse의 스파스
인덱스(sparse index)가 \texttt{WHERE venue='kalshi'} 조건을 인덱스 단계에서부터
효율적으로 잘라낼 수 있다는 부수적인 이득도 있다.
% src: 00.architecture.mvp1.md §3.4

\begin{example}[venue가 없었다면 벌어졌을 일]
\label{ex:venue-collision}
가상의 상황을 생각해 보자. Kalshi 티커 \texttt{KXBTCD-26SEP22-B112000}과 Polymarket의
어떤 CLOB token id가 (실제로는 형태가 전혀 달라 일어나지 않지만) 문자열로 우연히
같다고 가정하면, \texttt{venue} 없이 \texttt{ORDER BY (ticker, time,
idempotency\_key)}만 쓸 때 두 거래소의 서로 다른 오더북 행이 같은 \texttt{time},
같은 \texttt{idempotency\_key}를 우연히 공유하는 순간 하나로 접혀 버린다. 서로
다른 시장의 데이터가 뒤섞이는 것이다. \texttt{venue}를 선두에 추가하면
\texttt{(kalshi, KXBTCD-...)}와 \texttt{(polymarket, KXBTCD-...)}는 애초에 다른
정렬 키를 갖게 되어 이 위험이 사라진다.
\end{example}

\section{왜 Kafka 트랜잭션(EOS)이 아니라 DB 중복 제거인가}
\label{sec:delivery-why-db}

\paragraph{무슨 문제가 있었나.} Kafka는 exactly-once semantics(EOS)라는 자체
트랜잭션 기능을 제공한다. "그러면 그걸 쓰면 되지 않느냐"는 질문이 자연스럽게
나온다.

\paragraph{그래서 무엇을 바꿨나.} 쓰지 않기로 했다. 결정적인 이유는 하나다.
\term{ClickHouse의 \texttt{ENGINE=Kafka} 테이블은 트랜잭션 컨슈머(transactional
consumer)가 아니다.} Kafka EOS는 트랜잭션을 지원하는 컨슈머 위에서만 의미가 있는데,
이 파이프라인의 읽기 쪽(ClickHouse Kafka 엔진)이 애초에 그 조건을 만족하지 못한다.
% src: 00.architecture.mvp1.md §5.5

\paragraph{왜 그 변경이 문제를 푸는가.} DB 단 중복 제거는 Kafka EOS보다 \term{적용
범위가 더 넓다}. Kafka EOS는 Kafka 내부의 read-process-write 패턴만 커버하지만,
DB 단 중복 제거는 producer 크래시, 브로커 재시도, 컨슈머 재읽기까지 \term{경로에
상관없이} 전부 커버한다. 또한 producer가 재시작해 PID(프로세스 ID)가 바뀌어도
\texttt{idempotency\_key}가 안정적으로 재계산되는 한 중복 제거가 그대로 유지된다.
Kafka EOS는 PID가 바뀌면 트랜잭션 보장이 사라진다.
% src: 00.architecture.mvp1.md §5.5

\begin{table}[htbp]
  \centering
  \small
  \begin{tabularx}{\linewidth}{@{}lXX@{}}
    \toprule
     & DB 중복 제거 (채택) & Kafka 트랜잭션 EOS \\
    \midrule
    적용 범위 & producer 크래시, 브로커 재시도, 컨슈머 재읽기 전부 & Kafka 내부
      read-process-write만 \\
    ClickHouse Kafka 엔진 & 지원 & 트랜잭션 컨슈머 미지원. 애초에 쓸 수 없음 \\
    producer 재시작 후 & \texttt{idempotency\_key}가 안정적이면 중복 제거 유지 &
      PID가 바뀌면 보장 소멸 \\
    비용 & 병합 시 CPU, 읽기 시 \texttt{FINAL} & 처리량 저하, 운영 복잡도 \\
    \bottomrule
  \end{tabularx}
  \caption{DB 단 중복 제거와 Kafka 트랜잭션 EOS 비교.}
  \label{tab:delivery-eos-vs-dedup}
\end{table}
% src: 00.architecture.mvp1.md §5.5

\section{정직하게 짚는 잔존 위험}
\label{sec:delivery-honest}

\paragraph{무슨 문제가 있었나.} 지금까지의 논의를 "이 파이프라인은 exactly-once다"로
요약하면 과장이다. \term{effectively-exactly-once}라는 표현을 쓴 이유는, 그 보장이
성립하지 않는 지점이 실제로 남아 있기 때문이다. 이 절은 그 지점을 숨기지 않고
나열한다.

\paragraph{(1) 중복이 아니라 유실로 바뀌는 지점: \texttt{kafka\_skip\_broken\_messages
= 1}.} ClickHouse Kafka 엔진은 이 설정 때문에, 들어온 메시지가 proto 스키마와 맞지
않으면 그 행을 \term{그냥 버린다}. 이것은 지금까지 이야기한 "중복"과 성질이 다르다.
중복은 나중에 접히기라도 하지만, 이 경우는 \term{처음부터 유실}이다. 스키마를 바꿀
때(\Cref{ch:contract}의 3단계 계약)가 이 파이프라인 전체에서 가장 위험한 순간인
이유가 여기에 있다. proto, codec, DDL 세 단계의 순서를 지키는 규율이 이 위험을 막는
유일한 수단이다.
% src: 00.architecture.mvp1.md §5.3-4

\paragraph{(2) \texttt{acks=all}이 사실상 \texttt{acks=1}인 지점.}
\Cref{sec:delivery-atleastonce}에서 다룬 것처럼, replication factor가 1이라
\texttt{acks=all}의 실질 효과는 \texttt{acks=1}이다. 재시도는 있지만, 그 재시도가
기대는 브로커가 단 한 대뿐이라는 사실은 바뀌지 않는다. 그 한 대의 디스크가 완전히
망가지면, 재시도로도 메꿀 수 없는 유실이 난다.
% src: 00.architecture.mvp1.md §5.3-5

\paragraph{(3) 순서가 보장되지 않는 지점.} Kafka 토픽은 전부 파티션이 1개이므로
얼핏 순서가 유지될 것 같지만, 실제로는 그렇지 않다. producer가
\texttt{enable.idempotence: false}로 설정되어 있고, librdkafka의
\texttt{max.in.flight.requests.per.connection}(연결당 동시 전송 가능한 요청 수)
기본값이 작지 않아, \term{재시도된 배치가 원래보다 늦게 도착}할 수 있다. 파티션이
1개라는 사실은 이 문제를 막아 주지 않는다.
% src: 00.architecture.mvp1.md §5.3-6

ClickHouse의 이력 테이블은 \texttt{time} 기준으로 정렬해서 읽으므로 순서 뒤집힘이
무해하다. 문제는 Redis에 저장하는 \term{현재 상태}다. 오더북 라이터가 오래된
스냅샷을 나중에 받으면, 손대지 않을 경우 최신 상태를 오래된 값으로 덮어써 버린다.
그래서 Redis 쓰기에는 별도의 방어가 필요하다. \texttt{seq}(시퀀스 번호)가 아니라
\texttt{time\_ms}를 비교하는 \term{last-write-wins(LWW, 나중에 쓴 값이 이김)} 가드를
Lua 스크립트로 구현해, 기존에 저장된 \texttt{time\_ms}보다 작거나 같은 스냅샷은
무시한다. \texttt{seq}가 아니라 \texttt{time\_ms}를 쓰는 이유는 \texttt{seq}가
재접속 시 리셋될 수 있는 값이기 때문이다. 이 가드는 \texttt{ob:\{venue\}:\{ticker\}}
단위로 걸리므로, 한 거래소의 순서 뒤집힘이 다른 거래소의 상태에 영향을 주지 않는다.
% src: 00.architecture.mvp1.md §4.4, §5.3-6, §5.4

\paragraph{(4) 중복 제거는 최종적(eventual)이다.} ReplacingMergeTree의 병합은
백그라운드에서 비동기로 일어난다. 병합이 일어나기 전에는 같은 이벤트의 중복 행이
그대로 조회에 노출된다. 이 파이프라인은 BTC 1종의 데이터량에서 \texttt{SELECT ...
FROM t FINAL}(병합을 강제로 앞당겨 읽는 옵션)로 이 문제를 단순하게 해결한다. 대시보드의
"현재 오더북"은 애초에 ClickHouse가 아니라 Redis에서 읽으므로 이 지연의 영향을
받지 않는다.
% src: 00.architecture.mvp1.md §5.6

\section{한 줄로 보는 경로별 요약}
\label{sec:delivery-summary}

\Cref{tab:delivery-summary}는 이 장에서 다룬 네 구간을 한 표로 모은 것이다.

\begin{table}[htbp]
  \centering
  \small
  \begin{tabularx}{\linewidth}{@{}lXX@{}}
    \toprule
    구간 & 보장 수준 & 근거 \\
    \midrule
    producer 큐 $\to$ Kafka 브로커 & 보장 없음. 최대 2초분 조용히 유실 가능 &
      메모리 큐, \texttt{dr\_cb: false}, WS는 재생 불가 \\
    Kafka 브로커 $\to$ ClickHouse Kafka 엔진 & at-least-once &
      librdkafka 재시도, MV 삽입 후 offset 커밋 \\
    ClickHouse Kafka 엔진 $\to$ 데이터 테이블 & 중복은 발생하되 전부 바이트 동일 &
      재시도/재읽기가 같은 payload를 반복 \\
    데이터 테이블 병합 & effectively-exactly-once (최종적) &
      \texttt{ReplacingMergeTree} + 정렬 키 완전 일치 \\
    \bottomrule
  \end{tabularx}
  \caption{producer 큐부터 ClickHouse 데이터 테이블까지, 구간별 전달 보장 요약.}
  \label{tab:delivery-summary}
\end{table}

\begin{example}[숫자로 보는 세 가지 시나리오]
\label{ex:delivery-scenarios}
같은 Kalshi 오더북 이벤트 하나(\texttt{venue=kalshi},
\texttt{ticker=KXBTCD-26SEP22-B112000}, \texttt{time\_ms=1758500000123},
\texttt{seq=4711})를 예로 세 가지 경우를 비교한다.
% src: 00.architecture.mvp1.md §3.5 (예시 값 재사용)

\textbf{(a) 브로커 ack 유실 후 재시도, 그리고 삽입-커밋 사이 크래시.} 같은 바이트가
producer에서 두 번, ClickHouse에서 다시 두 번 재처리되어 데이터 테이블에 임시로
4개의 행이 쌓일 수 있다. 네 행 모두 \texttt{(venue, ticker, time,
idempotency\_key)}가 동일하므로, 병합 후에는 1개 행만 남는다.

\textbf{(b) producer 재시작 후 같은 오더북 상태를 새 시각으로 재발행.} 이것은
재시도가 아니라 새 이벤트 생성이다. 새 \texttt{time\_ms}를 얻으므로 정렬 키가
달라지고, 이력에는 2개의 서로 다른 행으로 남는다. 이것이
\Cref{sec:delivery-dedup}에서 말한 "불변식을 깨는 유일한 경우"의 실제 사례다.

\textbf{(c) 같은 티커 문자열을 가진 두 거래소.} \texttt{venue}가 다르므로
Redis 키는 \texttt{ob|kalshi|KXBTCD-26SEP22-B112000|4711}과
\texttt{ob|polymarket|KXBTCD-26SEP22-B112000|4711}로 처음부터 갈라지고, ClickHouse
정렬 키도 \texttt{venue} 컬럼이 다르므로 두 거래소의 행이 병합에서 하나로 접히는
일은 없다.
% src: 00.architecture.mvp1.md §4.4, §5.4
\end{example}

\paragraph{정리.} producer 메모리 큐부터 브로커까지는 at-least-once조차 아니다.
크래시가 나면 최대 2초분이 조용히 사라지고, \texttt{dr\_cb: false} 때문에 그 유실은
로그에도 남지 않는다. 브로커부터 ClickHouse 데이터 테이블까지는 at-least-once이며,
그 구간의 모든 중복은 바이트 단위로 동일한 재전송이므로 \texttt{idempotency\_key}를
정렬 키에 포함한 \texttt{ReplacingMergeTree}가 effectively-exactly-once로 접는다.
다만 순서는 보장되지 않으므로(idempotence off), Redis의 현재 상태는
\texttt{time\_ms} 기준 last-write-wins 가드로 별도로 방어한다.
% src: 00.architecture.mvp1.md §5.7
